\section{الأعمال ذات الصلة}
\label{sec:related_work}

استعرض \EN{Majeed} وآخرون (2021) \cite{math9212829} تقنيات إخفاء المعلومات النصية السابقة لعصر نماذج اللغة الكبيرة، أي قبل الحقبة الحالية للمحولات. ويشير \EN{Setiadi} وآخرون (2025) \cite{Setiadi_Ghosal_Sahu_2025} إلى أن نماذج اللغة الكبيرة قد أعادت تنشيط إخفاء المعلومات اللغوي، ودرسوا طرائق حديثة بين 2021 و2025 تستخدم \EN{GPT-2} \cite{radford2019gpt2} و\EN{GPT-3} \cite{brown2020language} و\EN{LLaMA2} \cite{touvron2023llama2} و\EN{Baichuan2} \cite{xiao2024baichuan2}. ومع ذلك، فإن هاتين الورقتين أقرب إلى مسوح أو مراجعات، ولا تجعلان التوافق السياقي أو الطرق المرتكزة إلى نماذج اللغة الكبيرة معيارًا مركزيًا، كما أنهما ليستا مراجعتين منهجيتين شاملتين للأدبيات. والأهم أن أيًا منهما لا يضع \emph{التوافق السياقي} في المقدمة بوصفه معيارًا أساسيًا للحكم على مدى ملاءمة النص الحامل المولد للخطاب والمهمة والإطار الحواري المحيط. والأمر الجوهري أن المجال تطور من ``تضمينات المتجهات الإحصائية''، مثل \EN{Word2Vec} و\EN{GloVe}، إلى ``تضمينات متجهات نموذج اللغة'' التي تستغل محولات بمستوى \EN{BERT} وفضاءات دلالية أعلى بعدًا.

ينتج عن ذلك فراغ منهجي: فلا توجد مراجعة منهجية ترسم بصورة شاملة كيف أعادت المحولات واسعة النطاق تعريف إخفاء المعلومات النصي. وتتجاوز التطورات الحديثة التوليد البسيط إلى أطر متقدمة للتوليد النصي القابل للتحكم \EN{Controllable Text Generation (CTG)} \cite{zhang2020linguistic}. وتستخدم هذه الأطر المشفرات التلقائية التباينية \EN{VAEs} لنمذجة الخصائص الكامنة، ونماذج الانتشار لإدخال العشوائية، بما يحد من الارتباطات الزائفة بين الأسرار وشروط التحكم. وقد ركزت المسوح الكلاسيكية على استبدال المرادفات، والتلاعب بالمسافات، وترميز \EN{Huffman} \cite{math9212829}، وهي تقنيات سبقت نماذج اللغة الكبيرة. واعتمدت الطرائق الأقدم على القواعد الخالية من السياق أو سلاسل ماركوف، ما كان ينتج غالبًا نصوصًا سليمة نحويًا لكنها غير متماسكة دلاليًا. أما الطرق المعاصرة فتستفيد من تعلم الموجهات وضبط السوابق، بما يتيح تخصيصًا فعالًا للنموذج دون كلفة الضبط الكامل.

وينبغي أن تتطور استراتيجيات الدفاع تبعًا لذلك. فالكشف التقليدي عن الإخفاء، المبني على سمات إحصائية مصممة يدويًا، يتعثر أمام التخفي الإحصائي العالي في الإخفاء التوليدي. وعلى البحث الحالي أن يواجه ``البرمجيات الخبيثة الإخفائية''، أي الهجمات التي تخفي اتصالات القيادة والتحكم داخل وسائط رقمية بريئة ظاهريًا.

يعالج الإطار المقترح هذه الفجوات معًا: فهو ينتقل من التعديل على مستوى الرموز إلى الترميز عبر قناة اختيار دلالية، ويدمج إعادة بناء السياق الحتمية لضمان المزامنة بين المرسل والمستقبل دون الاعتماد على وصول الصندوق الأبيض إلى النموذج، وهو ما يميزه عن طرائق الأخذ التوليدي للعينات السابقة.
